% GENERATED FILE. Edit canonical lessons, then run npm run generate:books.
% canonical-source-hash: fnv1a64:d100bc4b207b3fc4
% canonical-lessons: GE-C96-laden, GE-C96-markt, GE-C96-bahnhof, GE-C96-flugzeug, GE-C96-auto

\chapter{Shop, Market, Railway station, Aeroplane, Car}
\label{ch:ge-shop-market-railway-station-aeroplane-car}
\hlchaptermodality{96}

\begin{chapteropening}
\textbf{By the end of this chapter:} \emph{I can say a shop, a market, a railway station, an aeroplane and a car.}

Complete the last lesson of chapter 96: I can say a shop, a market, a railway station, an aeroplane and a car.
\end{chapteropening}

\section[der Laden]{der Laden — a shop}

\label{lesson:GE-C96-laden}



\begin{quote}
Before the new one: say the German for far, then the German for on the right.
\end{quote}

\subsection*{You'll want to know: der Laden}
\textbf{der Laden} — \textquotedblleft{}a shop\textquotedblright{}.

\begin{cousinweb}
\textbf{Laden} $\leftarrow$ Middle High German \emph{lade}, \textquotedblleft{}board, plank, shutter\textquotedblright{} $\leftarrow$ Old High German \emph{lado} $\leftarrow$ Proto-Germanic \emph{*laþô} (EtymDB 2.1) — inherited: a shop was first the shutter-board goods were laid out on.
\end{cousinweb}

\begin{grammarlens}[title={What you've built}]
One more word for this chapter.
\end{grammarlens}

\subsection*{Your turn}
\begin{itemize}
\raggedright
  \item \emph{Say it:} \emph{der Laden}
  \item \emph{Say it:} \emph{der Laden}, once more
  \item \emph{Say it:} say \emph{der Laden}
  \item \emph{Recall:} say the German for far, then the German for on the right, then say \emph{der Laden} again
\end{itemize}

\begin{tcolorbox}[breakable,colback=teal!4,colframe=teal!35!black,title={Before you move on}]
What does der Laden mean? (\textquotedblleft{}A shop\textquotedblright{}.) Say it once more.
\end{tcolorbox}
\section[der Markt]{der Markt — a market}

\label{lesson:GE-C96-markt}



\begin{quote}
Before the new one: say the German for on the right, then the German for a shop.
\end{quote}

\subsection*{You'll want to know: der Markt}
\textbf{der Markt} — \textquotedblleft{}a market\textquotedblright{}.

\begin{cousinweb}
\textbf{Markt} $\leftarrow$ Old High German \emph{marcāt} $\leftarrow$ Latin \emph{mercātus}, \textquotedblleft{}trade, market\textquotedblright{} (EtymDB 2.1) — borrowed, like English \emph{market}.
\end{cousinweb}

\begin{grammarlens}[title={What you've built}]
One more word for this chapter.
\end{grammarlens}

\subsection*{Your turn}
\begin{itemize}
\raggedright
  \item \emph{Say it:} \emph{der Markt}
  \item \emph{Say it:} \emph{der Markt}, once more
  \item \emph{Say it:} say \emph{der Markt}
  \item \emph{Recall:} say the German for on the right, then the German for a shop, then say \emph{der Markt} again
\end{itemize}

\begin{tcolorbox}[breakable,colback=teal!4,colframe=teal!35!black,title={Before you move on}]
What does der Markt mean? (\textquotedblleft{}A market\textquotedblright{}.) Say it once more.
\end{tcolorbox}
\section[der Bahnhof]{der Bahnhof — a railway station}

\label{lesson:GE-C96-bahnhof}



\begin{quote}
Before the new one: say the German for a shop, then the German for a market.
\end{quote}

\subsection*{You'll want to know: der Bahnhof}
\textbf{der Bahnhof} — \textquotedblleft{}a railway station\textquotedblright{}.

\begin{cousinweb}
\textbf{Bahnhof} is built from the word for \textquotedblleft{}way, track\textquotedblright{} (Old High German \emph{bana}, \textquotedblleft{}path\textquotedblright{}) and the word for \textquotedblleft{}yard, court\textquotedblright{} (Old High German \emph{hof}) — both inherited (EtymDB 2.1): the railway's yard.
\end{cousinweb}

\begin{grammarlens}[title={What you've built}]
One more word for this chapter.
\end{grammarlens}

\subsection*{Your turn}
\begin{itemize}
\raggedright
  \item \emph{Say it:} \emph{der Bahnhof}
  \item \emph{Say it:} \emph{der Bahnhof}, once more
  \item \emph{Say it:} say \emph{der Bahnhof}
  \item \emph{Recall:} say the German for a shop, then the German for a market, then say \emph{der Bahnhof} again
\end{itemize}

\begin{tcolorbox}[breakable,colback=teal!4,colframe=teal!35!black,title={Before you move on}]
What does der Bahnhof mean? (\textquotedblleft{}A railway station\textquotedblright{}.) Say it once more.
\end{tcolorbox}
\section[das Flugzeug]{das Flugzeug — an aeroplane}

\label{lesson:GE-C96-flugzeug}



\begin{quote}
Before the new one: say the German for a market, then the German for a railway station.
\end{quote}

\subsection*{You'll want to know: das Flugzeug}
\textbf{das Flugzeug} — \textquotedblleft{}an aeroplane\textquotedblright{}.

\begin{cousinweb}
\textbf{Flugzeug} is built from the word for \textquotedblleft{}flight\textquotedblright{}, Old High German \emph{flug} $\leftarrow$ Proto-Germanic \emph{*flugiz}, and the word for \textquotedblleft{}gear, stuff\textquotedblright{}, Old High German \emph{giziug} (EtymDB 2.1): a flying device.
\end{cousinweb}

\begin{grammarlens}[title={What you've built}]
One more word for this chapter.
\end{grammarlens}

\subsection*{Your turn}
\begin{itemize}
\raggedright
  \item \emph{Say it:} \emph{das Flugzeug}
  \item \emph{Say it:} \emph{das Flugzeug}, once more
  \item \emph{Say it:} say \emph{das Flugzeug}
  \item \emph{Recall:} say the German for a market, then the German for a railway station, then say \emph{das Flugzeug} again
\end{itemize}

\begin{tcolorbox}[breakable,colback=teal!4,colframe=teal!35!black,title={Before you move on}]
What does das Flugzeug mean? (\textquotedblleft{}An aeroplane\textquotedblright{}.) Say it once more.
\end{tcolorbox}
\section[das Auto]{das Auto — a car}

\label{lesson:GE-C96-auto}



\begin{quote}
Before the new one: say the German for a railway station, then the German for an aeroplane.
\end{quote}

\subsection*{You'll want to know: das Auto}
\textbf{das Auto} — \textquotedblleft{}a car\textquotedblright{}.

\begin{cousinweb}
\textbf{Auto} is short for the longer word for \textquotedblleft{}automobile\textquotedblright{}, taken from French and built from Greek \emph{autós}, \textquotedblleft{}self\textquotedblright{}, and Latin \emph{mōbilis}, \textquotedblleft{}movable\textquotedblright{}: a carriage that moves itself. EtymDB 2.1 has the French \emph{auto} as its source.
\end{cousinweb}

\begin{grammarlens}[title={What you've built}]
One more word for this chapter.
\end{grammarlens}

\subsection*{Your turn}
\begin{itemize}
\raggedright
  \item \emph{Say it:} \emph{das Auto}
  \item \emph{Say it:} \emph{das Auto}, once more
  \item \emph{Say it:} say \emph{das Auto}
  \item \emph{Recall:} say the German for a railway station, then the German for an aeroplane, then say \emph{das Auto} again
\end{itemize}

\begin{tcolorbox}[breakable,colback=teal!4,colframe=teal!35!black,title={Before you move on}]
What does das Auto mean? (\textquotedblleft{}A car\textquotedblright{}.) Say it once more.
\end{tcolorbox}
